\section{Introduction}


%Presents a clear learning context and includes:
%1) the general scientific concept with the necessary background;
%2) primary goals and the specific scientific concept to be investigated;
%2) specific theory, experimental design and equations needed with fully defined quantities;
%3) any assumptions made or required;
%4) hypothesis to be investigated;
%5) references to supporting sources.

Applying a magnetic field across a ferromagnetic material changes its magnetic response.
In the case of steel, this response is characterised by a hysteresis curve, where ramping up the current produces a
different path compared to reversing the current.
In this experiment we aim to demonstrate the B-H hysteresis curve of a steel toroid by measuring the voltage in the 
coil as the current is adjusted up and down, as well as the magnetic induction.

The magnetic field strength, denoted as $H$, is found using the voltage and resistance of the coil.

\begin{equation}
            \label{eq:H}
    |\vec{H}| = H = \frac{N \cdot i}{L} = \frac{N}{2 \pi r} \frac{V}{R}
\end{equation}

Where, $N$ is the number of turns in the toroid, $i$ is the current (also may be calculated using Ohm's Law), $R$ is
the resistance of the coil, and $
r$ is the radius.

The work done per volume of the material may be found as

\begin{equation}
    \frac{W}{V} = \oint \vec{B} \cdot  d\vec{H} \label{eq:work}
\end{equation}

Where $W$ is the work done, $V$ is the volume of the material, $\vec{B}$ is the magnetic induction, and $d\vec{H}$ 
is the magnetic field strength element.

Magnetic susceptibility, $\chi$, is the quantity which describes how well the material aligns with the applied magnetic 
field, or how easily it is magnetised.
In ferromagnetic materials, this can be calculated using 

\begin{equation}
    B = \mu_0(1 + \chi) H\label{eq:chi}
\end{equation}

Where $\mu_0$ is the permeability of free space.

%Ref equation: Eqs.~\eqref{eq:norm_damp}, ~\eqref{eq:crit_damp}, and \eqref{eq:osc_damp}, respectively.